\frfilename{mt271.tex}
\versiondate{11.12.08}
\copyrightdate{1995}

\def\chaptername{Teori probabilitas}
\def\sectionname{Distribusi}

\newsection{271}

Saya memulai bab ini dengan pembahasan mengenai `distribusi
probabilitas', yaitu ukuran probabilitas pada $\BbbR^n$ yang didefinisikan
oleh keluarga $(X_1,\ldots,X_n)$ dari variabel-variabel acak. Saya memberikan
hasil-hasil dasar yang menguraikan keadaan ketika dua distribusi sama
(271G), integrasi terhadap suatu distribusi (271E), dan fungsi kerapatan
probabilitas (271H-271K).   %271H 271I 271J 271K

\medskip

\leader{271A}{Notasi }\cmmnt{Saya baru saja menghabiskan beberapa paragraf
untuk mencoba menguraikan perbedaan hakiki antara teori probabilitas dan
teori ukuran. Namun ada uji yang lebih cepat untuk mengetahui apakah penulis
Anda seorang ahli teori ukuran atau ahli probabilitas: bukalah sembarang
halaman, lalu cari frasa `fungsi terukur' dan `variabel acak', serta rumus
`$\int fd\mu$' dan `$\Expn(X)$'. Anggota pertama dari setiap pasangan akan
memungkinkan Anda mendiagnosis `ukuran', dan anggota kedua `probabilitas',
dengan kemungkinan keliru yang kecil. Sejauh ini dalam risalah ini saya
dengan teguh menggunakan istilah para ahli teori ukuran, dengan beberapa
kekhasan pribadi. Namun dalam sebuah bab tentang teori probabilitas, saya
mendapati bahwa notasi teori ukuran, meskipun sepenuhnya memadai dalam arti
formal, begitu merusak rumusan-rumusan yang sudah akrab sehingga membuatnya
tidak wajar. Lagi pula, pada suatu saat -- jika belum -- Anda tentu harus
mengenal bahasa para ahli probabilitas. Jadi dalam bab ini saya akan mengambil
langkah besar ke arah itu. Untungnya, saya pikir hal ini dapat dilakukan tanpa
menimbulkan pertentangan langsung, sehingga dalam jilid-jilid selanjutnya saya
akan dapat menggunakan pekerjaan ini dalam notasi mana pun yang pada saat itu
tampak sesuai, tanpa perlu merumuskannya kembali.

\medskip

}{\bf (a)}\dvro{ Misalkan}{ Jadi misalkan} $(\Omega,\Sigma,\mu)$ adalah
suatu ruang probabilitas. \cmmnt{Saya menggunakan kesempatan yang diberikan
oleh suatu frasa baru untuk mengambil langkah teknis.} Suatu
{\bf variabel acak bernilai real}
pada $\Omega$ akan merupakan anggota dari
$\eusm L^0(\mu)$\cmmnt{, sebagaimana didefinisikan dalam 241A; yaitu, suatu
fungsi bernilai real $X$ yang didefinisikan pada suatu himpunan bagian
koterabaikan dari $\Omega$ sedemikian sehingga $X$ terukur terhadap
pelengkapan $\hat\mu$ dari $\mu$, atau, jika Anda lebih suka,
sedemikian sehingga $X\restr E$ terukur terhadap $\Sigma$ untuk suatu
himpunan koterabaikan $E\subseteq\Omega$}.\cmmnt{\footnote{Untuk uraian
tentang bagaimana istilah ini menjadi baku, lihat
{\tt
http://www.dartmouth.edu/$\sptilde$chance/\penalty-100Doob/\penalty-100conversation.html}.}}
% http://www.dartmouth.edu/~chance/Doob/conversation.html  % Apr09

\header{271Ab}{\bf (b)} Jika $X$ adalah variabel acak bernilai real pada
suatu ruang probabilitas
$(\Omega,\Sigma,\mu)$, tuliskan $\Expn(X)=\int X\,d\mu$ jika nilai ini
terdefinisi dalam $[-\infty,\infty]$\cmmnt{ dalam arti Bab 12 dan \S133}.
Dalam hal ini saya akan menyebut $\Expn(X)$ sebagai {\bf rataan} atau
{\bf ekspektasi} dari $X$. Jadi kita boleh mengatakan bahwa
`$X$ {\bf mempunyai ekspektasi hingga}' sebagai ganti `$X$
terintegralkan'. \cmmnt{133A mengatakan bahwa
`$\Expn(X+Y)=\Expn(X)+\Expn(Y)$ setiap kali $\Expn(X)$
dan $\Expn(Y)$ beserta jumlahnya terdefinisi dalam $[-\infty,\infty]$', dan
122P menjadi `suatu variabel acak bernilai real $X$ mempunyai ekspektasi
hingga jika dan hanya jika $\Expn(|X|)<\infty$'.}

\header{271Ac}{\bf (c)} Jika $X$ adalah variabel acak bernilai real dengan
ekspektasi hingga, {\bf varians} dari $X$ adalah

\Centerline{$\Var(X)=\Expn(X-\Expn(X))^2
=\Expn(X^2-2\Expn(X)X+\Expn(X)^2)=\Expn(X^2)-(\Expn(X))^2$\dvro{.}{}}

\cmmnt{\noindent (Perhatikan bahwa rumus ini menunjukkan bahwa
$\Expn(X)^2\le\Expn(X^2)$; bandingkan 244Xd(i).) $\Var(X)$ berhingga jika
dan hanya jika $\Expn(X^2)<\infty$, yaitu, jika dan hanya jika
$X\in\eusm L^2(\mu)$ (244A). Secara khusus, $X+Y$ dan $cX$ mempunyai
varians hingga setiap kali $X$ dan $Y$ demikian dan $c\in\Bbb R$.}


\header{271Ad}{\bf (d)} Saya akan membolehkan diri menggunakan rumus-rumus
seperti

\Centerline{$\Pr(X>a)$,\quad $\Pr(X-\epsilon\le Y\le X+\delta)$,}

\noindent dengan $X$ dan $Y$ variabel acak pada ruang probabilitas yang sama
$(\Omega,\Sigma,\mu)$, untuk masing-masing berarti

\Centerline{$\hat\mu\{\omega:\omega\in\dom X,\,X(\omega)>a\}$,}

\Centerline{$\hat\mu\{\omega:\omega\in\dom X\cap\dom Y,\,
X(\omega)-\epsilon\le Y(\omega)\le X(\omega)+\delta\}$,}

\noindent dengan menuliskan $\hat\mu$ untuk pelengkapan $\mu$\cmmnt{
seperti biasa}. \cmmnt{Ada dua hal yang perlu diperhatikan di sini.
Pertama, $\Pr$ bergantung pada $\hat\mu$, bukan pada $\mu$; pada dasarnya,
notasi tersebut secara otomatis mengarahkan kita untuk melengkapi ruang
probabilitas $(\Omega,\Sigma,\mu)$. Tentu saja, saya dapat pula menulis

\Centerline{$\Pr(X^2+Y^2>1)
=\mu^*\{\omega:\omega\in\dom X\cap\dom Y,\,
X(\omega)^2+Y(\omega)^2>1\}$,}

\noindent dengan mengambil $\mu^*$ sebagai ukuran luar pada $\Omega$ yang
berkaitan dengan $\mu$ (132B). Kedua,} saya akan menggunakan notasi ini
{\it hanya untuk predikat yang bersesuaian dengan himpunan terukur Borel};
artinya, saya akan menulis

\Centerline{$\Pr(\psi(X_1,\ldots,X_n))
=\hat\mu\{\omega:\omega\in\bigcap_{i\le n}\dom X_i,\,
\psi(X_1(\omega),\ldots,X_n(\omega))\}$}

\noindent hanya ketika himpunan

\Centerline{$\{(\alpha_1,\ldots,\alpha_n):
\psi(\alpha_1,\ldots,\alpha_n)\}$}

\noindent merupakan himpunan Borel dalam $\BbbR^n$. \cmmnt{Sebagian alasan
bagi pembatasan ini akan tampak dalam beberapa paragraf berikutnya;
$\Pr(\psi(X_1,\ldots,X_n))$ harus merupakan sesuatu yang dapat dihitung dari
pengetahuan mengenai distribusi gabungan $X_1,\ldots,X_n$, sebagaimana
didefinisikan dalam 271C. Sebenarnya kita dapat dengan aman memperluas gagasan
ini ke predikat $\psi$ yang `terukur universal', yang akan dibahas dalam Jilid
4. Namun dapat saja terjadi bahwa $\mu$ memberi ukuran kepada himpunan
berbentuk $\{\omega:X(\omega)\in A\}$ untuk suatu himpunan $A$ yang luar biasa
tak beraturan, dan dalam kasus semacam itu akan lebih bijaksana memandang hal
ini sebagai patologi kebetulan dari ruang probabilitas, lalu menanganinya
dengan cara yang agak berbeda.

(Saya melihat bahwa saya telah agak gegabah menganggap bahwa rumus di atas
mendefinisikan $\Pr(\psi(X_1,\ldots,X_n))$ untuk setiap predikat Borel $\psi$.
Ini merupakan konsekuensi dari 271Bb di bawah.)
}

\leader{271B}{Teorema} Misalkan $(\Omega,\Sigma,\mu)$ suatu ruang
probabilitas, dan $X_1,\ldots,X_n$ variabel-variabel acak bernilai real pada
$\Omega$. Tetapkan
$\pmb{X}(\omega)=(X_1(\omega),\ldots,X_n(\omega))$ untuk
$\omega\in\bigcap_{i\le n}\dom X_i$.

(a) Terdapat tepat satu ukuran Radon $\nu$ pada $\BbbR^n$ sedemikian sehingga

\Centerline{$\nu\ocint{-\infty,a}=\Pr(X_i\le\alpha_i$
untuk setiap $i\le n)$}

\noindent setiap kali $a=(\alpha_1,\ldots,\alpha_n)
\in\BbbR^n$\cmmnt{, dengan menuliskan $\ocint{-\infty,a}$ untuk
$\prod_{i\le n}\ocint{-\infty,\alpha_i}$};

(b) $\nu\BbbR^n=1$ dan $\nu E=\hat\mu(\pmb{X}^{-1}[E])$ setiap kali $\nu E$
terdefinisi, dengan $\hat\mu$ pelengkapan dari $\mu$; khususnya,
$\nu E=\Pr((X_1,\ldots,X_n)\in E)$ untuk setiap himpunan Borel
$E\subseteq\BbbR^n$.

\proof{ Misalkan $\hat\Sigma$ domain dari $\hat\mu$, dan tetapkan
$D=\bigcap_{i\le n}\dom X_i=\dom\pmb{X}$; maka $D$ koterabaikan, sehingga
termasuk dalam $\hat\Sigma$. Misalkan
$\hat\mu_D=\hat\mu\restrp\Cal PD$ ukuran subruang pada $D$ (131B, 214B), dan
$\nu_0$ ukuran citra $\hat\mu_D\pmb{X}^{-1}$ (234D); misalkan $\Tau$ domain
dari $\nu_0$.

Tuliskan $\Cal B$ untuk aljabar himpunan Borel dalam $\BbbR^n$. Maka
$\Cal B\subseteq\Tau$. \Prf\ Untuk $i\le n$, $\alpha\in\Bbb R$ tetapkan
$F_{i\alpha}=\{x:x\in\BbbR^n,\,\xi_i\le\alpha\}$,
$H_{i\alpha}=\{\omega:\omega\in\dom X_i,\,X_i(\omega)\le\alpha\}$.
$X_i$ terukur terhadap $\hat\Sigma$ dan domainnya berada dalam
$\hat\Sigma$, sehingga $H_{i\alpha}\in\hat\Sigma$, dan
$\pmb{X}^{-1}[F_{i\alpha}]=D\cap H_{i\alpha}$ terukur terhadap
$\hat\mu_D$. Jadi $F_{i\alpha}\in\Tau$.
Karena $\Tau$ adalah aljabar-sigma bagian-bagian dari $\BbbR^n$,
$\Cal B\subseteq\Tau$ (121J).\ \Qed

Dengan demikian $\nu_0\restr\Cal B$ adalah ukuran pada $\BbbR^n$ dengan
domain $\Cal B$; tentu saja $\nu_0\BbbR^n=\hat\mu D=1$. Berdasarkan 256C,
pelengkapan $\nu$ dari $\nu_0\restr\Cal B$ adalah ukuran Radon pada
$\BbbR^n$, dan $\nu\BbbR^n=\nu_0\BbbR^n=1$.

Untuk $E\in\Cal B$,

\Centerline{$\nu E=\nu_0E=\hat\mu_D\pmb{X}^{-1}[E]
=\hat\mu\pmb{X}^{-1}[E]
=\Pr((X_1,\ldots,X_n)\in E)$.}

\noindent Secara lebih umum, jika $E\in\dom\nu$, maka terdapat himpunan Borel
$E'$, $E''$ sedemikian sehingga $E'\subseteq E\subseteq E''$ dan
$\nu(E''\setminus E')=0$, sehingga
$\pmb{X}^{-1}[E']\subseteq\pmb{X}^{-1}[E]\subseteq\pmb{X}^{-1}[E'']$ dan
$\hat\mu(\pmb{X}^{-1}[E'']\setminus\pmb{X}^{-1}[E'])=0$. Ini berarti bahwa
$\pmb{X}^{-1}[E]\in\hat\Sigma$ dan

\Centerline{$\hat\mu\pmb{X}^{-1}[E]=\hat\mu\pmb{X}^{-1}[E']=\nu E'=\nu
E$.}

Adapun mengenai ketunggalan $\nu$, jika $\nuprime$ adalah suatu ukuran Radon
pada $\Bbb
R^n$ sedemikian sehingga
$\nuprime\ocint{-\infty,a}=\Pr(X_i\le\alpha_i\Forall i\le n)$ untuk setiap
$a\in\BbbR^n$, maka tentu

\Centerline{$\nuprime\BbbR^n
=\lim_{k\to\infty}\nuprime\ocint{-\infty,k\tbf{1}}
=\lim_{k\to\infty}\nu\ocint{-\infty,k\tbf{1}}=1=\nu\BbbR^n$.}

\noindent Selain itu, $\Cal I=\{\ocint{-\infty,a}:a\in\BbbR^n\}$ tertutup
terhadap irisan berhingga, dan $\nu$ serta $\nuprime$ bersepakat pada
$\Cal I$. Berdasarkan Teorema Kelas Monoton (atau lebih tepatnya, akibatnya
136C), $\nu$ dan $\nuprime$ bersepakat pada aljabar-sigma yang dibangkitkan
oleh $\Cal I$, yaitu $\Cal B$ (121J), dan keduanya identik (256D).
}%end of proof of 271B

\leader{271C}{Definisi} Misalkan $(\Omega,\Sigma,\mu)$ suatu ruang
probabilitas dan $X_1,\ldots,X_n$ variabel-variabel acak bernilai real pada
$\Omega$. Yang saya maksud dengan {\bf distribusi} ({\bf gabungan}) atau
{\bf hukum} $\nu_{\pmb{X}}$ dari keluarga
$\pmb{X}=(X_1,\ldots,X_n)$ ialah ukuran probabilitas Radon $\nu$ dalam 271B.
\cmmnt{Jika kita memandang $\pmb{X}$ sebagai fungsi dari
$\bigcap_{i\le n}\dom X_i$ ke $\BbbR^n$, maka
$\nu_{\pmb{X}}E=\Pr(\pmb{X}\in E)$ untuk setiap himpunan Borel
$E\subseteq\BbbR^n$.}

\leader{271D}{Catatan }\cmmnt{{\bf (a)} Pemilihan ukuran probabilitas Radon
$\nu_{\pmb{X}}$ sebagai `distribusi' dari $\pmb{X}$, dengan tuntutan agar
`ukuran Radon' lengkap, tentu agak sebarang. Selain prinsip umum bahwa ukuran
selalu sebaiknya dilengkapi, konvensi-konvensi ini lebih sesuai dengan
sebagian pekerjaan dalam Jilid 4 dan dengan hasil-hasil seperti 272G di bawah.

\spheader 271Db
Perhatikan bahwa agar dapat berbicara tentang distribusi suatu keluarga
$\pmb{X}=(X_1,\ldots,X_n)$ dari variabel acak, semua $X_i$ harus didefinisikan
pada ruang probabilitas yang sama.

\spheader 271Dc Saya melihat bahwa bahasa yang telah saya pilih
memungkinkan $X_i$ mempunyai domain yang berbeda-beda, sehingga keluarga
$(X_1,\ldots,X_n)$ mungkin tidak dapat diidentifikasi secara tepat dengan
fungsi yang bersesuaian dari $\bigcap_{i\le n}\dom X_i$ ke $\BbbR^n$.
Namun saya berharap penggunaan lambang yang sama $\pmb{X}$ untuk keduanya
tidak menimbulkan kebingungan.

\spheader 271Dd Tidaklah berguna memandang seluruh ukuran citra
$\nu_0=\hat\mu_D\pmb{X}^{-1}$ dalam bukti 271B sebagai distribusi dari
$\pmb{X}$, kecuali jika kebetulan ukuran itu sama dengan
$\nu=\nu_{\pmb{X}}$. `Distribusi' suatu variabel acak adalah tepat aspek
variabel tersebut yang dapat dipisahkan dari segala pertimbangan mengenai ruang
yang mendasari $(\Omega,\Sigma,\mu)$, dan pokok hasil seperti 271K dan 272G
ialah bahwa distribusi dapat dihitung dari satu sama lain, tanpa kembali ke
model variabel acak yang relatif mudah berubah dan tak pasti sebagai suatu
fungsi pada ruang probabilitas.

\medskip

}{\bf (e)} %271De
Jika $\pmb{X}=(X_1,\ldots,X_n)$ dan
$\pmb{Y}=(Y_1,\ldots,Y_n)$ sedemikian sehingga $X_i\eae Y_i$ untuk setiap
$i$, maka\cmmnt{

\Centerline{$\{\omega:\omega\in\bigcap_{i\le n}\dom X_i,
  \,X_i(\omega)\le\alpha_i\Forall i\le n\}
\symmdiff\{\omega:\omega\in\bigcap_{i\le n}\dom Y_i,
  \,Y_i(\omega)\le\alpha_i\Forall i\le n\}$}

\noindent terabaikan, sehingga

$$\eqalign{\Pr(X_i\le\alpha_i\Forall i\le n)
&=\hat\mu\{\omega:\omega\in\bigcap_{i\le n}\dom X_i,
  \,X_i(\omega)\le\alpha_i\Forall i\le n\}\cr
&=\Pr(Y_i\le\alpha_i\Forall i\le n)\cr}$$

\noindent untuk semua $\alpha_1,\ldots,\alpha_n\in\Bbb R$,
dan} $\nu_{\pmb{X}}=\nu_{\pmb{Y}}$. \cmmnt{Ini berarti bahwa, jika kita
mau, kita dapat memandang suatu distribusi sebagai ukuran
$\nu_{\pmb{u}}$ dengan $\pmb{u}=(u_0,\ldots,u_n)$ barisan berhingga dalam
$L^0(\mu)$. Dalam bab ini saya tidak akan menekankan pendekatan tersebut,
tetapi pendekatan itu akan selalu ada di benak saya.
}%end of comment

\leader{271E}{Fungsi terukur dari variabel acak: Proposisi}
Misalkan $\pmb{X}=(X_1,\ldots,X_n)$ suatu keluarga variabel acak
\cmmnt{(seperti biasa dalam konteks semacam itu, yang saya maksud ialah bahwa
semuanya berada pada ruang probabilitas yang sama
$(\Omega,\Sigma,\mu)$)}; tuliskan $\Tau_{\pmb{X}}$ untuk domain distribusi
$\nu_{\pmb{X}}$, dan misalkan $h$ suatu fungsi bernilai real yang terukur
terhadap $\Tau_{\pmb{X}}$ dan didefinisikan hampir di mana-mana terhadap
$\nu_{\pmb{X}}$ pada $\BbbR^n$. Maka kita mempunyai variabel acak
$Y=h(X_1,\ldots,X_n)$ yang didefinisikan dengan menetapkan

\Centerline{$h(X_1,\ldots,X_n)(\omega)
=h(X_1(\omega),\ldots,X_n(\omega))$ untuk setiap
$\omega\in\pmb{X}^{-1}[\dom h]$.}

\noindent Distribusi $\nu_Y$ dari $Y$ adalah ukuran pada $\Bbb R$ yang
didefinisikan oleh rumus

\Centerline{$\nu_YF=\nu_{\pmb{X}}h^{-1}[F]$}

\noindent untuk dan hanya untuk himpunan-himpunan $F\subseteq\Bbb R$ sedemikian sehingga
$h^{-1}[F]\in\Tau_{\pmb{X}}$. Selain itu

\Centerline{$\Expn(Y)=\int h\,d\nu_{\pmb{X}}$}

\noindent dalam arti bahwa jika salah satu dari keduanya ada dalam
$[-\infty,\infty]$, maka yang lain juga ada dan keduanya sama.

\proof{{\bf (a)(i)} Sekali lagi, tuliskan
$(\Omega,\hat\Sigma,\hat\mu)$ untuk pelengkapan
$(\Omega,\Sigma,\mu)$. Karena

\Centerline{$\Omega\setminus\dom Y
\subseteq\bigcup_{i\le n}(\Omega\setminus\dom X_i)
  \cup\pmb{X}^{-1}[\BbbR^n\setminus\dom h]$}

\noindent terabaikan (dengan menggunakan 271Bb), $\dom Y$ koterabaikan.
Jika $a\in\Bbb R$, maka

\Centerline{$E=\{x:x\in\dom h,\,h(x)\le a\}\in\Tau_{\pmb{X}},$}

\noindent sehingga

\Centerline{$\{\omega:\omega\in\Omega,\,Y(\omega)\le a\}
=\pmb{X}^{-1}[E]\in\hat\Sigma$.}

\noindent Karena $a$ sebarang, $Y$ terukur terhadap $\hat\Sigma$, dan
merupakan variabel acak.

\medskip

\quad{\bf (ii)} Misalkan $\tilde h:\BbbR^n\to\Bbb R$ suatu perluasan dari
$h$ ke seluruh $\BbbR^n$. Maka $\tilde h$ terukur terhadap
$\Tau_{\pmb{X}}$, sehingga ukuran citra biasa
$\nu_{\pmb{X}}\tilde h^{-1}$, yang didefinisikan pada
$\{F:\tilde h^{-1}[F]\in\dom\nu_{\pmb{X}}\}$, merupakan ukuran probabilitas
Radon pada $\Bbb R$ (256G). Namun untuk setiap $A\subseteq\Bbb R$,

\Centerline{$\tilde h^{-1}[A]\symmdiff h^{-1}[A]
\subseteq\BbbR^n\setminus\dom h$}

\noindent terabaikan terhadap $\nu_{\pmb{X}}$, sehingga
$\nu_{\pmb{X}}h^{-1}[F]=\nu_{\pmb{X}}\tilde h^{-1}[F]$ jika salah satunya
terdefinisi.

Jika $F\subseteq\Bbb R$ suatu himpunan Borel, maka

\Centerline{$\nu_YF=\hat\mu\{\omega:Y(\omega)\in F\}
=\hat\mu(\pmb{X}^{-1}[h^{-1}[F]])=\nu_{\pmb{X}}(h^{-1}[F])$.}

\noindent Jadi $\nu_Y$ dan $\nu_{\pmb{X}}\tilde h^{-1}$ bersepakat pada
himpunan-himpunan Borel dan sama (256D lagi).

\medskip

{\bf (b)} Sekarang terapkan Teorema 235E pada ukuran $\hat\mu$ dan
$\nu_{\pmb{X}}$ serta fungsi $\phi=\pmb{X}$. Kita mempunyai

\Centerline{$\int\chi(\pmb{X}^{-1}[F])d\hat\mu
=\hat\mu(\pmb{X}^{-1}[F])=\nu_{\pmb{X}}F$}

\noindent untuk setiap $F\in\Tau_{\pmb{X}}$, berdasarkan 271Bb. Karena $h$
terukur secara virtual terhadap $\nu_{\pmb{X}}$ dan didefinisikan hampir di
mana-mana terhadap $\nu_{\pmb{X}}$, 235Eb memberi tahu kita bahwa

\Centerline{$\int h(\pmb{X})d\mu=\int h(\pmb{X})d\hat\mu
=\int h\,d\nu_{\pmb{X}}$}

\noindent setiap kali salah satu ruas terdefinisi dalam
$[-\infty,\infty]$, dan inilah tepat hasil yang kita perlukan.
}%end of proof of 271E

\vleader{72pt}{271F}{Akibat} Jika $X$ suatu variabel acak tunggal dengan
distribusi $\nu_{X}$, maka

\Centerline{$\Expn(X)=\int_{-\infty}^{\infty}x\,\nu_X(dx)$}

\noindent jika salah satunya terdefinisi dalam $[-\infty,\infty]$.
Demikian pula

\Centerline{$\Expn(X^2)=\int_{-\infty}^{\infty} x^2\,\nu_X(dx)$}

\noindent (apa pun $X$). Jika $X$, $Y$ dua variabel acak\cmmnt{ (pada ruang
probabilitas yang sama!)} maka kita mempunyai

\Centerline{$\Expn(X\times Y)=\int xy\,\nu_{(X,Y)}d(x,y)$}

\noindent jika salah satu ruas terdefinisi dalam $[-\infty,\infty]$.

\medskip

\noindent{\bf Catatan}\dvAnew{2013} Jika $\nu$ distribusi suatu variabel acak
bernilai real, yaitu suatu ukuran probabilitas Radon pada $\Bbb R$, saya akan
mengatakan bahwa {\bf ekspektasi} $\Expn(\nu)$ dari $\nu$ adalah
$\int_{-\infty}^{\infty}x\,\nu(dx)$ jika nilai ini terdefinisi; jika $\nu$
mempunyai ekspektasi hingga, maka {\bf varians} $\Var(\nu)$ adalah
$\int x^2\,\nu(dx)-(\Expn(\nu))^2$. Jadi jika $X$ suatu variabel acak bernilai
real dengan distribusi $\nu_X$, $\Expn(X)=\Expn(\nu_X)$ dan
$\Var(X)=\Var(\nu_X)$ setiap kali nilai-nilai ini terdefinisi.

\leader{271G}{Fungsi distribusi (a)} Jika $X$ suatu variabel acak bernilai
real, {\bf fungsi distribusi}-nya adalah fungsi $F_X:\Bbb R\to[0,1]$ yang
didefinisikan dengan menetapkan

\Centerline{$F_X(a)=\Pr(X\le a)=\nu_X\ocint{-\infty,a}$}

\noindent untuk setiap $a\in\Bbb R$. ({\bf Peringatan!} beberapa penulis
lebih menyukai $F_X(a)=\Pr(X<a)$.) \cmmnt{Perhatikan bahwa $F_X$ tak
menurun, bahwa $\lim_{a\to-\infty}F_X(a)=0$, bahwa
$\lim_{a\to\infty}F_X(a)=1$ dan bahwa
$\lim_{x\downarrow a}F_X(x)=F_X(a)$ untuk setiap $a\in\Bbb R$.
Berdasarkan 271Ba,} $X$ dan $Y$ mempunyai distribusi yang sama jika dan hanya
jika $F_X=F_Y$.

\spheader 271Gb
Jika $X_1,\ldots,X_n$ variabel-variabel acak bernilai real pada ruang
probabilitas yang sama, {\bf fungsi distribusi (gabungan)} mereka adalah
fungsi $F_{\pmb{X}}:\BbbR^n\to[0,1]$ yang didefinisikan dengan menuliskan

\Centerline{$F_{\pmb{X}}(a)=\Pr(X_i\le\alpha_i\Forall i\le n)$}

\noindent setiap kali $a=(\alpha_1,\ldots,\alpha_n)\in\BbbR^n$. Jika
$\pmb{X}$ dan $\pmb{Y}$ mempunyai fungsi distribusi yang sama, keduanya
mempunyai distribusi yang sama\cmmnt{, berdasarkan versi berdimensi $n$ dari
271B}.

\leader{271H}{Kerapatan} Misalkan $\pmb{X}=(X_1,\ldots,X_n)$ suatu keluarga
variabel acak, semuanya didefinisikan pada ruang probabilitas yang sama. Suatu
{\bf fungsi kerapatan} untuk $(X_1,\ldots,X_n)$ adalah turunan
Radon-Nikod\'ym, terhadap ukuran Lebesgue, bagi distribusi
$\nu_{\pmb{X}}$\cmmnt{; yaitu, suatu fungsi nonnegatif $f$, yang
terintegralkan terhadap ukuran Lebesgue $\mu_L$ pada $\BbbR^n$, sedemikian
sehingga

\Centerline{$\int_Efd\mu_L=\nu_{\pmb{X}}E=\Pr(\pmb{X}\in E)$}

\noindent untuk setiap himpunan Borel $E\subseteq\BbbR^n$ (256J) -- tentu
saja jika fungsi semacam itu ada}.

\leader{271I}{Proposisi} Misalkan $\pmb{X}=(X_1,\ldots,X_n)$ suatu keluarga
variabel acak, semuanya didefinisikan pada ruang probabilitas yang sama.
Tuliskan $\mu_L$ untuk ukuran Lebesgue pada $\BbbR^n$.

(a) Terdapat fungsi kerapatan untuk $\pmb{X}$ jika dan hanya jika
$\Pr(\pmb{X}\in E)=0$ untuk setiap himpunan Borel $E$ sedemikian sehingga
$\mu_LE=0$.

(b) Suatu fungsi nonnegatif $f$ yang terintegralkan secara Lebesgue adalah
fungsi kerapatan untuk $\pmb{X}$ jika dan hanya jika
$\int_{\ocint{-\infty,a}}fd\mu_L=\Pr(\pmb{X}\in\ocint{-\infty,a})$ untuk
setiap $a\in\BbbR^n$.

(c) Misalkan $f$ suatu fungsi kerapatan untuk $\pmb{X}$, dan
$G=\{x:f(x)>0\}$.
Maka jika $h$ suatu fungsi bernilai real terukur Lebesgue yang didefinisikan
hampir di mana-mana dalam $G$,

\Centerline{$\Expn(h(\pmb{X}))=\int h\,d\nu_{\pmb{X}}
=\int h\times fd\mu_L$}

\noindent jika salah satu dari ketiga integral itu terdefinisi dalam
$[-\infty,\infty]$, dengan menafsirkan $(h\times f)(x)$ sebagai $0$ jika
$f(x)=0$ dan $x\notin\dom h$.

\proof{{\bf (a)} Terapkan 256J pada ukuran probabilitas Radon
$\nu_{\pmb{X}}$.

\medskip

{\bf (b)} Tentu saja syarat tersebut perlu. Jika syarat itu dipenuhi, maka
(berdasarkan teorema B.Levi)

\Centerline{$\int fd\mu_L
=\lim_{k\to\infty}\int_{\ocint{-\infty,k\pmb{1}}}fd\mu_L
=\lim_{k\to\infty}\nu_{\pmb{X}}\ocint{-\infty,k\pmb{1}}
=1$.}

\noindent Jadi kita mempunyai ukuran probabilitas Radon $\nu$ yang
didefinisikan dengan menuliskan

\Centerline{$\nu E=\int_Efd\mu_L$}

\noindent setiap kali $E\cap\{x:f(x)>0\}$ terukur Lebesgue (256E). Kita
mengandaikan bahwa
$\nu\ocint{-\infty,a}=\nu_{\pmb{X}}\ocint{-\infty,a}$ untuk setiap
$a\in\BbbR^n$; berdasarkan 271Ba, seperti biasa, $\nu=\nu_{\pmb{X}}$,
sehingga

\Centerline{$\int_Efd\mu_L=\nu E=\nu_{\pmb{X}}E
=\Pr(\pmb{X}\in E)$}

\noindent untuk setiap himpunan Borel $E\subseteq\BbbR^n$, dan $f$ adalah
fungsi kerapatan untuk $\pmb{X}$.

\medskip

{\bf (c)} Berdasarkan 256E, $\nu_{\pmb{X}}$ adalah ukuran integral tak tentu
di atas $\mu_L$ yang berkaitan dengan $f$. Jadi, dengan menuliskan
$G=\{x:f(x)>0\}$, kita mempunyai

\Centerline{$\int h\,d\nu_{\pmb{X}}=\int h\times fd\mu_L$}

\noindent setiap kali salah satunya terdefinisi dalam $[-\infty,\infty]$
(235K). Berdasarkan 234La, $h$ terukur terhadap $\Tau_{\pmb{X}}$ dan
didefinisikan hampir di mana-mana terhadap $\nu_{\pmb{X}}$, dengan
$\Tau_{\pmb{X}}=\dom\nu_{\pmb{X}}$, sehingga
$\Expn(h(\pmb{X}))=\int h\,d\nu_{\pmb{X}}$ berdasarkan 271E.
}%end of proof of 271I

\leader{271J}{}\cmmnt{ Perangkat yang dikembangkan dalam \S263 memadai
untuk memberikan hasil yang sangat umum mengenai kerapatan variabel acak
berbentuk $\phi(\pmb{X})$, sebagai berikut.

\medskip

\noindent}{\bf Teorema} Misalkan $\pmb{X}=(X_1,\ldots,X_n)$ suatu keluarga
variabel acak, dan $D\subseteq\BbbR^n$ suatu himpunan Borel sedemikian sehingga
$\Pr(\pmb{X}\in D)=1$. Misalkan $\phi:D\to\BbbR^n$ suatu fungsi yang dapat
didiferensialkan relatif terhadap domainnya di setiap titik dalam $D$; untuk
$x\in D$, misalkan $T(x)$ suatu turunan dari $\phi$ pada $x$, dan tetapkan
$J(x)=|\det T(x)|$. Andaikan bahwa $J(x)\ne 0$ untuk setiap $x\in D$, bahwa
$\pmb{X}$ mempunyai fungsi kerapatan $f$; dan andaikan pula bahwa
$\sequence{k}{D_k}$ suatu barisan himpunan Borel yang saling lepas,
yang gabungannya $D$, sedemikian sehingga $\phi_k=\phi\restr D_k$ injektif
untuk setiap $k$. Maka $\phi(\pmb{X})$ mempunyai fungsi kerapatan
$g=\sum_{k=0}^{\infty}g_k$ dengan

$$\eqalignno{g_k(y)&=\Bover{f(\phi_k^{-1}(y))}{J(\phi_k^{-1}(y))}
\text{ untuk }y\in\phi[D_k\cap\dom f],\cr
&=0\text{ untuk }y\in\BbbR^n\setminus\phi[D_k].\cr}$$

\proof{ Berdasarkan 262Ia, $\phi$ kontinu, sehingga terukur Borel, dan karena
itu $\phi(\pmb{X})$ merupakan variabel acak.

Untuk sementara, tetapkan $k\in\Bbb N$ dan suatu himpunan Borel
$F\subseteq\BbbR^n$. Berdasarkan 263D(iii), $\phi[D_k]$ terukur, dan
berdasarkan 263D(ii) $\phi[D_k\setminus\dom f]$ terabaikan. Fungsi $g_k$
sedemikian sehingga $f(x)=J(x)g_k(\phi(x))$ untuk setiap
$x\in D_k\cap\dom f$, sehingga berdasarkan 263D(v) kita mempunyai

$$\eqalign{\int_Fg_k\,d\mu
&=\int_{\phi[D_k]} g_k\times\chi F\,d\mu
=\int_{D_k} J(x)g_k(\phi(x))\chi F(\phi(x))\mu(dx)\cr
&=\int_{D_k\cap\phi^{-1}[F]}fd\mu
=\Pr(\pmb{X}\in D_k\cap\phi^{-1}[F]).\cr}$$

\noindent (Integral
$\int_{\phi[D_k]} g_k\times\chi F$ terdefinisi karena
$\int_{D_k}J\times(g_k\times\chi F)\phi$ terdefinisi, dan integral
$\int
g_k\times\chi F$ terdefinisi karena $\phi[D_k]$ terukur dan $g_k$ nol di luar
$\phi[D_k]$.)

Sekarang jumlahkan atas $k$. Setiap $g_k$ nonnegatif, sehingga berdasarkan
teorema B.Levi (123A, 123Xa)

$$\eqalign{\int_Fg\,d\mu
&=\sum_{k=0}^{\infty}\int_Fg_k\,d\mu
=\sum_{k=0}^{\infty}\Pr(\pmb{X}\in D_k\cap\phi^{-1}[F])\cr
&=\Pr(\pmb{X}\in\phi^{-1}[F])
=\Pr(\phi(\pmb{X})\in F).\cr}$$

\noindent Karena $F$ sebarang, $g$ adalah fungsi kerapatan untuk
$\phi(\pmb{X})$, sebagaimana diklaim.
}%end of proof of 271J

\leader{271K}{}\cmmnt{ Penerapan teorema terakhir pada transformasi biasa
kadang-kadang tidak langsung, sehingga saya memberikan sebuah contoh.

\medskip

\noindent}{\bf Proposisi} Misalkan $X$, $Y$ dua variabel acak dengan suatu
fungsi kerapatan gabungan $f$. Maka $X\times Y$ mempunyai fungsi kerapatan
$h$, dengan

\Centerline{$h(u)
=\int_{-\infty}^{\infty}\Bover1{|v|}f(\bover{u}{v},v)dv$}

\noindent setiap kali nilai ini terdefinisi dalam $\Bbb R$.

\proof{ Tetapkan $\phi(x,y)=(xy,y)$ untuk $x$, $y\in\BbbR^2$. Maka $\phi$
dapat didiferensialkan, dengan turunan $T(x,y)=\Matrix{y&x\\0&1}$, sehingga
$J(x,y)=|\det T(x,y)|=|y|$. Tetapkan $D=\{(x,y):y\ne 0\}$; maka $D$ adalah
himpunan Borel koterabaikan dalam $\BbbR^2$ dan $\phi\restr D$ injektif.
Sekarang $\phi[D]=D$ dan $\phi^{-1}(u,v)=(\bover{u}{v},v)$ untuk $v\ne 0$.
Jadi $\phi(X,Y)=(X\times Y,Y)$ mempunyai fungsi kerapatan $g$, dengan

\Centerline{$g(u,v)=\Bover{f(u/v,v)}{|v|}$ jika $v\ne 0$.}

Untuk menemukan fungsi kerapatan bagi $X\times Y$, kita menghitung

\Centerline{$\Pr(X\times Y\le a)
=\int_{\ocint{-\infty,a}\times\Bbb R}g
=\int_{-\infty}^a\int_{-\infty}^{\infty}g(u,v)dv\,du
=\int_{-\infty}^ah$}

\noindent berdasarkan teorema Fubini (252B, 252C). Khususnya, $h$ terdefinisi
dan berhingga hampir di mana-mana; dan berdasarkan 271Ib, fungsi tersebut
adalah fungsi
kerapatan untuk $X\times Y$.
}%end of proof of 271K

\leader{*271L}{}\cmmnt{ Ketika suatu variabel acak disajikan sebagai limit
suatu barisan variabel acak, hasil berikut dapat sangat berguna.

\medskip

\noindent}{\bf Proposisi} Misalkan $\sequencen{X_n}$ suatu barisan variabel
acak bernilai real yang konvergen dalam ukuran menuju suatu variabel acak
$X$\cmmnt{ (definisi: 245A)}. Dengan menuliskan $F_{X_n}$, $F_X$ untuk
fungsi distribusi masing-masing $X_n$, $X$,

\Centerline{$F_X(a)=\inf_{b>a}\liminf_{n\to\infty}F_{X_n}(b)
=\inf_{b>a}\limsup_{n\to\infty}F_{X_n}(b)$}

\noindent untuk setiap $a\in\Bbb R$.

\proof{ Tetapkan $\gamma=\inf_{b>a}\liminf_{n\to\infty}F_{X_n}(b)$,
$\gamma'=\inf_{b>a}\limsup_{n\to\infty}F_{X_n}(b)$.

\medskip

{\bf (a)} $F_X(a)\le\gamma$. \Prf\ Ambil sembarang $b>a$ dan $\epsilon>0$.
Maka terdapat $n_0\in\Bbb N$ sedemikian sehingga
$\Pr(|X_n-X|\ge b-a)\le\epsilon$ untuk setiap $n\ge n_0$ (245F). Sekarang,
untuk $n\ge n_0$,

\Centerline{$F_X(a)
=\Pr(X\le a)
\le\Pr(X_n\le b)+\Pr(X_n-X\ge b-a)
\le F_{X_n}(b)+\epsilon$.}

\noindent Jadi $F_X(a)\le\liminf_{n\to\infty}F_{X_n}(b)+\epsilon$; karena
$\epsilon$ sebarang, $F_X(a)\le\liminf_{n\to\infty}F_{X_n}(b)$; karena
$b$ sebarang, $F_X(a)\le\gamma$.\ \Qed

\medskip

{\bf (b)} $\gamma'\le F_X(a)$. \Prf\ Misalkan $\epsilon>0$. Maka terdapat
$\delta>0$ sedemikian sehingga $F_X(a+2\delta)\le F_X(a)+\epsilon$ (271Ga).
Selanjutnya, terdapat $n_0\in\Bbb N$ sedemikian sehingga
$\Pr(|X_n-X|\ge\delta)\le\epsilon$ untuk setiap $n\ge n_0$.
Dalam hal ini, untuk $n\ge n_0$,

$$\eqalign{F_{X_n}(a+\delta)
&=\Pr(X_n\le a+\delta)
\le\Pr(X\le a+2\delta)+\Pr(X-X_n\ge\delta)\cr
&\le F_X(a+2\delta)+\epsilon
\le F_X(a)+2\epsilon.\cr}$$

\noindent Dengan demikian

\Centerline{$\gamma'\le\limsup_{n\to\infty}F_{X_n}(a+\delta)
\le F_X(a)+2\epsilon$.}

\noindent Karena $\epsilon$ sebarang, $\gamma'\le F_X(a)$.\ \Qed

\medskip

{\bf (c)} Karena tentu saja $\gamma\le\gamma'$, haruslah
$F_X(a)=\gamma=\gamma'$, sebagaimana diklaim.
}%end of proof of 271L

\exercises{
\leader{271X}{Latihan dasar $\pmb{>}$(a)}
%\spheader 271Xa
Misalkan $X$ suatu variabel acak bernilai real dengan ekspektasi hingga, dan
$\epsilon>0$. Tunjukkan bahwa
$\Pr(|X-\Expn(X)|\ge\epsilon)\le\Bover1{\epsilon^2}\Var(X)$.
(Ini adalah {\bf ketaksamaan Chebyshev}.)
%271A

\sqheader 271Xb Misalkan $F:\Bbb R\to[0,1]$ suatu fungsi tak menurun
sedemikian sehingga
(i) $\lim_{a\to-\infty}F(a)=0$ (ii) $\lim_{a\to\infty}F(a)=1$ (iii)
$\lim_{x\downarrow a}F(x)=F(a)$ untuk setiap $a\in\Bbb R$. Tunjukkan bahwa
terdapat tepat satu ukuran probabilitas Radon $\nu$ pada $\Bbb R$ sedemikian
sehingga $F(a)=\nu\ocint{-\infty,a}$ untuk setiap
$a\in\Bbb R$. \Hint{114Xa.} Selanjutnya tunjukkan bahwa $F$
adalah fungsi distribusi dari suatu variabel acak.
%271G

\sqheader 271Xc Misalkan $X$ suatu variabel acak bernilai real dengan fungsi
kerapatan $f$. (i) Tunjukkan bahwa $|X|$ mempunyai fungsi kerapatan $g_1$
dengan $g_1(x)=f(x)+f(-x)$ setiap kali $x\ge 0$ dan $f(x)$, $f(-x)$ keduanya
terdefinisi, dan $0$ jika tidak. (ii) Tunjukkan bahwa $X^2$ mempunyai fungsi
kerapatan $g_2$ dengan
$g_2(x)=(f(\sqrt{x})+f(-\sqrt{x}))/2\sqrt{x}$ setiap kali $x>0$ dan nilai ini
terdefinisi, serta $0$ untuk $x$ lainnya. (iii) Tunjukkan bahwa jika
$\Pr(X=0)=0$ maka $1/X$ mempunyai fungsi kerapatan $g_3$ dengan
$g_3(x)=\bover1{x^2}f(\bover1x)$ setiap kali nilai ini terdefinisi.
(iv) Tunjukkan bahwa jika $\Pr(X<0)=0$ maka $\sqrt{X}$ mempunyai fungsi
kerapatan $g_4$ dengan $g_4(x)=2xf(x^2)$ jika $x\ge 0$ dan $f(x^2)$
terdefinisi, serta $0$ jika tidak.
%271J

\sqheader 271Xd Misalkan $X$ dan $Y$ variabel-variabel acak dengan fungsi
kerapatan gabungan $f:\BbbR^2\to\Bbb R$. Tunjukkan bahwa $X+Y$ mempunyai
fungsi kerapatan $h$ dengan $h(u)=\int f(u-v,v)dv$ untuk hampir setiap $u$.
%271K

\spheader 271Xe Misalkan $X$, $Y$ variabel-variabel acak dengan fungsi kerapatan
gabungan $f:\BbbR^2\to\Bbb R$. Tunjukkan bahwa $X/Y$ mempunyai fungsi
kerapatan $h$ dengan $h(u)=\int|v|f(uv,v)dv$ untuk hampir setiap $u$.
%271K

\spheader 271Xf Rancang suatu bukti alternatif untuk 271K dengan menggunakan
teorema Fubini dan substitusi satu dimensi untuk menunjukkan bahwa

\Centerline{$\int_a^b\int_{-\infty}^{\infty}
\Bover1{|v|}f(\bover{u}{v},v)dv\,du
=\int_{\{(u,v):a\le uv\le b\}}f$}

\noindent setiap kali $a\le b$ dalam $\Bbb R$.
%271K

\leader{271Y}{Latihan lanjutan (a)}
%\spheader 271Ya
Misalkan $\frak T$ topologi dari
$\BbbR^{\Bbb N}$ dan $\Cal B$ aljabar-sigma himpunan Borel (256Yf).
(i) Misalkan $\Cal I$ keluarga himpunan berbentuk

\Centerline{$\{x:x\in\BbbR^{\Bbb N},\,
x(i)\le \alpha_i\Forall i\le n\}$,}

\noindent dengan $n\in\Bbb N$ dan $\alpha_i\in\Bbb R$
untuk setiap $i\le n$. Tunjukkan bahwa $\Cal B$ adalah keluarga terkecil
dari bagian-bagian $\BbbR^{\Bbb N}$ sedemikian sehingga ($\alpha$)
$\Cal I\subseteq\Cal B$
($\beta$) $B\setminus A\in\Cal B$ setiap kali $A$, $B\in\Cal B$ dan
$A\subseteq B$ ($\gamma$) $\bigcup_{k\in\Bbb N}A_k\in\Cal B$ untuk setiap
barisan tak menurun $\sequence{k}{A_k}$ dalam $\Cal B$. (ii) Tunjukkan bahwa
jika $\mu$, $\mu'$ dua ukuran berhingga total yang didefinisikan pada
$\BbbR^{\Bbb N}$, dan
$\mu F$ serta $\mu' F$ terdefinisi dan sama untuk setiap $F\in\Cal I$, maka
$\mu E$ dan $\mu' E$ terdefinisi dan sama untuk setiap $E\in\Cal B$.
(iii) Tunjukkan bahwa jika $\Omega$ suatu himpunan dan $\Sigma$ suatu
aljabar-sigma bagian-bagian dari $\Omega$ dan
$X:\Omega\to\BbbR^{\Bbb N}$ suatu fungsi, maka $X^{-1}[E]\in\Sigma$ untuk
setiap $E\in\Cal B$ jika dan hanya jika $\pi_iX$ terukur terhadap $\Sigma$
untuk setiap $i\in \Bbb N$, dengan $\pi_i(x)=x(i)$ untuk setiap
$x\in\BbbR^{\Bbb N}$, $i\in \Bbb N$. (iv) Tunjukkan bahwa jika
$\pmb{X}=\sequence{i}{X_i}$ suatu barisan variabel acak bernilai real pada ruang
probabilitas $(\Omega,\Sigma,\mu)$, maka terdapat tepat satu ukuran
probabilitas $\nu_{\pmb{X}}^{\Cal B}$, dengan domain $\Cal B$, sedemikian
sehingga
$\nu_{\pmb{X}}^{\Cal B}\{x:x(i)\le\alpha_i\Forall i\le n\}
=\Pr(X_i\le\alpha_i\Forall i\le n)$ untuk setiap
$\alpha_0,\ldots,\alpha_n\in\Bbb R$. (v) Di bawah syarat-syarat (iv),
tunjukkan bahwa terdapat tepat satu ukuran Radon $\nu_{\pmb{X}}$ pada
$\BbbR^{\Bbb N}$ (dalam arti 256Yf) sedemikian sehingga
$\nu_{\pmb{X}}\{x:x(i)\le\alpha_i\Forall i\le n\}
=\Pr(X_i\le\alpha_i\Forall i\le n)$ untuk setiap
$\alpha_0,\ldots,\alpha_n\in\Bbb R$.

\header{271Yb}{\bf (b)} Misalkan $F:\BbbR^2\to[0,1]$ suatu fungsi.
Tunjukkan bahwa pernyataan berikut ekuivalen: (i) $F$ adalah fungsi distribusi
dari suatu pasangan $(X_1,X_2)$ variabel acak (ii) terdapat suatu ukuran
probabilitas $\nu$ pada $\BbbR^2$ sedemikian sehingga
$\nu\ocint{-\infty,a}=F(a)$ untuk setiap $a\in\BbbR^2$
(iii)($\alpha$)
$F(\alpha_1,\alpha_2)+F(\beta_1,\beta_2)
\ge F(\alpha_1,\beta_2)+F(\beta_1,\alpha_2)$ setiap kali
$\alpha_1\le\beta_1$
dan $\alpha_2\le\beta_2$ ($\beta$) $F(\alpha_1,\alpha_2)
=\lim_{\xi_1\downarrow\alpha_1,\xi_2\downarrow\alpha_2}F(\xi_1,\xi_2)$
untuk setiap $\alpha_1$, $\alpha_2$ ($\gamma$)
$\lim_{\alpha\to-\infty}F(\alpha,\beta)
=\lim_{\alpha\to-\infty}F(\beta,\alpha)=0$ untuk semua $\beta$
($\delta$) $\lim_{\alpha\to\infty}F(\alpha,\alpha)=1$.
\Hint{untuk interval setengah terbuka tak kosong $\ocint{a,b}$, tetapkan
$\lambda\ocint{a,b}=F(\alpha_1,\alpha_2)+F(\beta_1,\beta_2)
  -F(\alpha_1,\beta_2)-F(\beta_1,\alpha_2)$, dan
lanjutkan seperti dalam 115B-115F.}

\header{271Yc}{\bf (c)} Perumum (b) ke dimensi yang lebih tinggi, dengan
mencari rumus yang sesuai untuk menggantikan rumus dalam (iii-$\alpha$) dari
(b).

\spheader 271Yd Misalkan $(\Omega,\Sigma,\mu)$ suatu ruang probabilitas dan
$\Cal F$ suatu filter pada $\eusm L^0(\mu)$ yang konvergen menuju
$X_0\in\eusm L^0(\mu)$ untuk topologi konvergensi dalam ukuran. Tunjukkan
bahwa, dengan menuliskan $F_X$ untuk fungsi distribusi dari
$X\in\eusm L^0(\mu)$,

\Centerline{$F_{X_0}(a)=\inf_{b>a}\liminf_{X\to\Cal F}F_X(b)
=\inf_{b>a}\limsup_{X\to\Cal F}F_X(b)$}

\noindent untuk setiap $a\in\Bbb R$.
%271L

\spheader 271Ye\dvAnew{2009.} Misalkan $X$, $Y$ variabel-variabel acak
nonnegatif dengan distribusi yang sama, dan
$h:\coint{0,\infty}\to\coint{0,\infty}$ suatu fungsi tak menurun. Tunjukkan
bahwa $\Expn(X\times hY)\le\Expn(Y\times hY)$.
\Hint{dalam bahasa 252Yo, $(Y\times hY)^*=Y^*\times(hY)^*$.}
}%end of exercises

\endnotes{
\Notesheader{271} Sebagian besar bagian ini tampaknya telah dipenuhi oleh
hal-hal teknis. Ini mungkin tidak mengherankan mengingat bagian ini membahas
hubungan antara suatu variabel acak vektor $\pmb{X}$ dan distribusi terkait
$\nu_{\pmb{X}}$, dan hal ini dengan sendirinya membawa kita ke ladang ranjau
yang telah saya coba petakan dalam \S235. Sesungguhnya, saya menggunakan
hasil dari \S235 dua kali; sekali dalam 271E, dengan
$\phi(\omega)=\pmb{X}(\omega)$ dan $J(\omega)=1$, dan sekali dalam 271I,
dengan $\phi(x)=x$ dan $J(x)=f(x)$.

Fungsi distribusi variabel acak satu dimensi mudah dicirikan (271Xb); dalam
dimensi yang lebih tinggi kita harus bekerja lebih keras (271Yb-271Yc).
Distribusi, alih-alih fungsi distribusi, dapat diuraikan untuk barisan tak
hingga variabel acak (271Ya); bahkan, gagasan ini dapat diperluas ke keluarga
tak terhitung, tetapi hal ini memerlukan teori ukuran topologis yang
sesungguhnya dan termasuk dalam Jilid 4.

Pernyataan 271J rumit, bahkan dapat dikatakan berbelit-belit. Pokoknya ialah
bahwa banyak transformasi $\phi$ yang paling penting tidak dengan sendirinya
injektif, tetapi dapat dengan mudah diuraikan menjadi fragmen-fragmen injektif
(lihat, misalnya, 271Xc dan 263Xd). Pokok 271K ialah bahwa kita sering ingin
menerapkan gagasan di sini pada transformasi yang singular, dan bahkan
mengubah dimensi variabel acak. Saya belum memberikan teorema yang membuat
penerapan semacam itu rutin, dan lebih menyarankan agar Anda mencari siasat
seperti yang digunakan dalam bukti 271K, yang bagaimanapun diperlukan jika
Anda menginginkan rumus-rumus yang mudah ditangani. Tentu saja metode lain
tersedia (271Xf).
}%end of notes

\discrpage
